\documentclass{article}
\usepackage{amsmath,amssymb}
\usepackage{xurl}
\usepackage[hidelinks]{hyperref}
% Shared commands for notes in docs/paper-gaps/.

\DeclareUnicodeCharacter{2080}{\ifmmode{}_{0}\else\textsubscript{0}\fi}
\DeclareUnicodeCharacter{2081}{\ifmmode{}_{1}\else\textsubscript{1}\fi}
\DeclareUnicodeCharacter{2082}{\ifmmode{}_{2}\else\textsubscript{2}\fi}
\DeclareUnicodeCharacter{2099}{\ifmmode{}_{n}\else\textsubscript{n}\fi}

\newcommand{\Lean}{\textsc{Lean}}
\ExplSyntaxOn
\NewDocumentCommand{\leanid}{m}{
  \ifmmode
    \textsf{\detokenize{#1}}
  \else
    \group_begin:
    \urlstyle{sf}
    \tl_set:Nn \l_tmpa_tl {#1}
    \tl_replace_all:Nnn \l_tmpa_tl {\_} {_}
    \exp_args:NV \path \l_tmpa_tl
    \group_end:
  \fi
}
\ExplSyntaxOff
\providecommand{\tr}{\operatorname{tr}}
\newcommand{\tnleanIssueBase}{https://github.com/LionSR/TNLean/issues/}
\newcommand{\tnleanPrBase}{https://github.com/LionSR/TNLean/pull/}
\newcommand{\tnleanDocBase}{https://sirui-lu.com/TNLean/docs/}
\newcommand{\ghissue}[1]{\href{\tnleanIssueBase#1}{GitHub issue \##1}}
\newcommand{\ghpr}[1]{\href{\tnleanPrBase#1}{PR \##1}}
\newcommand{\doclink}[1]{\href{\tnleanDocBase#1.html}{\path{#1}}}

% Dirac notation and the identity operator, as in blueprint/src/macros/common.tex.
% \providecommand keeps note-local definitions authoritative.
\usepackage{dsfont}
\providecommand{\ket}[1]{|#1\rangle}
\providecommand{\bra}[1]{\langle #1|}
\providecommand{\braket}[2]{\langle #1 | #2 \rangle}
\providecommand{\ketbra}[2]{|#1\rangle\!\langle #2|}
\providecommand{\Id}{\mathds{1}}
\providecommand{\one}{\mathds{1}}
\providecommand{\C}{\mathbb{C}}
\providecommand{\MN}[1]{M_{#1}(\C)}
\providecommand{\MD}{M_D(\C)}

% External referencing for paper sources.  The build helper writes sanitized
% auxiliary files under build/paper-aux/ and excludes duplicated source labels.
\usepackage{xr}

% Import a generated paper auxiliary file with a caller-supplied label prefix.
% The candidate paths support builds from docs/paper-gaps/, the repository root,
% and build/paper-aux/ itself.
\newcommand{\importpaperaux}[2]{%
  \IfFileExists{../../build/paper-aux/#2.aux}{%
    \externaldocument[#1]{../../build/paper-aux/#2}%
  }{%
    \IfFileExists{build/paper-aux/#2.aux}{%
      \externaldocument[#1]{build/paper-aux/#2}%
    }{%
      \IfFileExists{#2.aux}{%
        \externaldocument[#1]{#2}%
      }{}%
    }%
  }%
}

% General external reference macros
\newcommand{\paperref}[2]{\ref{#1-#2}}
\newcommand{\papereqref}[2]{(\ref{#1-#2})}

% CPSV16 source (1606.00608 / MPDO-22-12-17-2.tex) external document and shortcuts
\importpaperaux{cpsv16-}{1606.00608/MPDO-22-12-17-2-tnlean}
\newcommand{\cpsvref}[1]{\paperref{cpsv16}{#1}}
\newcommand{\cpsveqref}[1]{\papereqref{cpsv16}{#1}}

% Verdict marker: \gapnote{<kind>}{<status>}. Kind is the policy
% classification (clarification, local-correction, scope-restriction,
% unfaithful, false-source, open-gap); status is open, wip (elimination
% underway), resolved, or historical. Machine-read by the site index;
% prints a verdict line.
\newcommand{\gapnote}[2]{\par\noindent\textbf{Verdict:} #1 (#2).\par}

\title{The Depth Lower Bound in the Gauge of a Normal Tensor}
\author{}
\date{2026-09-26}
\begin{document}
\maketitle
\gapnote{scope-restriction}{open}

\section{What the source states}
Let $\ket{\phi_N}$ be the normalized translation-invariant matrix product
states on $N$ sites generated by a normal tensor $A$ with correlation length
$\xi=-1/\ln\lvert\lambda_2\rvert>0$, and let $\ket{\psi_N}$ be obtained from
product states by local circuits of depth $T$. Theorem~1
of~\cite{Malz2024Preparation} states that if $T=o(\log N)$, then
$\varepsilon(\phi_N,\psi_N)=1-\lvert\braket{\phi_N}{\psi_N}\rvert>1/2$ for all
$N$ larger than some $N_0$.\footnote{Source traceability: label \path{th:1}
    of the arXiv LaTeX source, restated in the Supplemental Material, section
    ``Proof of Theorem~1''; a copy of the source is kept at
    \path{References/2307.01696/main.tex}.} The source's definition of a normal
tensor includes the normalization $\lambda_1=1$ of the leading eigenvalue of
the transfer map $\mathcal E_A$, and the source remarks that after a gauge
transformation $\mathcal E_A=\ket{\rho}\bra{\mathbb 1}+R$, that is,
\begin{align}
    \sum_iA^{i\dagger}A^i & =\mathbb 1, &
    \mathcal E_A(\rho)    & =\rho,      &
    \rho                  & >0,         &
    \operatorname{Tr}\rho & =1.
    \label{eq:mswc24_depth_lower_bound_gauge}
\end{align}

\section{What is formalized}
The formalized statement assumes (\ref{eq:mswc24_depth_lower_bound_gauge}),
a length $L\ge1$ at which the products of $L$ matrices of $A$ span the full
matrix algebra, and an eigenvalue $\lambda_2\ne1$ of $\mathcal E_A$ whose
modulus bounds those of all other eigenvalues different from $1$, with
$-1/\ln\lvert\lambda_2\rvert>0$. Under these hypotheses it proves the source's
conclusion, and the quantitative form: there are $B,C,b>0$ such that
$\lvert\braket{\phi_N}{\psi}\rvert\ge1/2$ for a unit vector $\psi$ prepared in
depth $T$ and $N\ge B(T+1)$ forces $N\le C(T+1)e^{b(T+1)}$.\footnote{Formal
    declarations:
    \leanid{MPSTensor.eventually_one_half_lt_infidelity_of_isLittleO_log}
    and \leanid{MPSTensor.exists_depth_lower_bound} in
    \doclink{TNLean/MPS/Preparation/DepthLowerBound}.}

The restriction is the gauge (\ref{eq:mswc24_depth_lower_bound_gauge}). Both
the source's proof and the chapter's proof begin by passing to it: a normal
tensor becomes $\lambda_1^{-1/2}\sigma^{1/2}A^i\sigma^{-1/2}$ with $\sigma$ the
positive fixed point of the dual transfer map, which multiplies
$\ket{\phi_N(A)}$ by a positive constant and leaves the normalized vector and
the correlation length unchanged. That reduction is not formalized.

\section{Elimination plan}
Formalize the reduction to the gauge for a normal tensor: existence of the
positive fixed points of $\mathcal E_A$ and its dual, invariance of the
normalized vectors $\ket{\phi_N}$ under the rescaled conjugation, and
invariance of the spectrum of $\mathcal E_A$ up to the factor $\lambda_1$.
Then derive the source's Theorem~1 for every normal tensor from the gauge form.

\section{Proof route}
The formal proof follows the chapter's argument rather than the source's. The
source bounds a fidelity of blocked product states and uses its Lemma~2 at the
single separation $s=2T+1$ with vanishing one-point functions; the correction
of that lemma is recorded in
\path{docs/paper-gaps/mswc24_decaying_correlations_windowed_connected.tex}. The
chapter instead averages the centred observables over $n\approx N/\Delta$
windows spaced $\Delta=s'+L+2T$ apart and compares the averages in
$\ket{\phi_N}$ and $\ket{\psi_N}$ through the inequality
$\lvert\braket{\phi}{\psi}\rvert\,\lvert\langle Z\rangle_\phi-\langle
    Z\rangle_\psi\rvert\le\sigma_\phi(Z)+\sigma_\psi(Z)$.

\bibliographystyle{alpha}
\bibliography{references}
\end{document}
